\documentclass[conference,letterpaper]{IEEEtran}
\IEEEoverridecommandlockouts
\usepackage[T1]{fontenc}
\usepackage{lmodern}
\renewcommand{\ttdefault}{lmtt}
\usepackage{booktabs}
\usepackage[hyphens]{url}
\usepackage[hidelinks]{hyperref}

\begin{document}

\title{Compute-Hour Efficient Coverage Growth:\\
Baselines, Adaptive CHIA, and Parallel Resources}

\author{\IEEEauthorblockN{Khushal Sethi and Kapil Bansal}
\IEEEauthorblockA{CHIA Hackathon at A$^3$@MICRO 2026\\
Artifact: \url{https://github.com/kbansal1512/chia-hackathon}}}

\maketitle

\begin{abstract}
Verification campaigns share scarce compute across constrained-random
verification (CRV), directed collateral, and formal analysis. We frame the
hackathon artifact as a coverage-per-compute-hour comparison with three arms:
(B1)~CRV only; (B2)~a fixed CRV$\rightarrow$directed$+$formal hybrid that uses
the same engines as CHIA; and (B3)~a CHIA adaptive allocator on that same pool.
A CHIA systems claim requires B3 to beat B2, not merely B1. A second feature
targets scale: today's loop mostly occupies one resource class at a time;
with $N$ heterogeneous seats, a large cover space, and a large formal
partition, CHIA must admit and replan work under scarcity. We also ship a
source-pinned BOOM monitor package with partial MediumBOOM qualification
evidence. The monitors elaborate on generated MediumBoomV3Config RTL and pass
five instrumented workloads. The aggregate reaches 73.59\% full-design signal
toggle activity and 582 of 1,192 named cover counters. A DTLB cover witness
replays, a targeted ROB mutation is caught, and the generated integer and FP
RenameBusyTable assertion cones pass 2-inductive proofs. We do not claim CPU
closure: per-instance audit is partial, one FP activation instance remains
open, and a DTLB stale-refill counterexample needs owner review. Bounded cover
non-hits do not retire coverage.
\end{abstract}

\begin{IEEEkeywords}
CHIA, coverage per compute-hour, adaptive scheduling, formal scope, BOOM
\end{IEEEkeywords}

\section{Problem}

CRV, directed tests, and formal solvers are often available together. The hard
decision is which engine receives the next compute-hour---and, at scale, which
of $N$ seats receives the next job---while preserving evidence that progress
means what it appears to mean. Showing that formal reaches points simulation
cannot establishes engine complementarity. It does not by itself establish that
an adaptive control plane is better than a competent fixed hybrid schedule.

\section{Feature 1: three baselines}

We compare three arms. \textbf{B1 (CRV only)} is the diagnostic floor.
\textbf{B2 (static hybrid)} runs a fixed CRV$\rightarrow$directed$\rightarrow$formal
schedule, frozen before held-out outcomes, and uses the same engines as CHIA;
it is the primary control. \textbf{B3 (CHIA adaptive)} uses the same resource
pool but admits the next job from open debt and measured yield.

Primary outcomes are sound coverage closed per compute-hour and wall time (or
cost) to a pre-registered target. B3 must be compared to B2. Beating B1 alone
is insufficient because B1 withholds engines known to be necessary for some
obligations. The repository's MediumBOOM smoke run is pilot evidence of
complementarity (simulation bins versus Load/Store Unit cover properties). It
is not a published B1/B2/B3 ROI table.

\section{Feature 2: one seat today, \texorpdfstring{$N$}{N} seats next}

The present driver largely serializes engine classes: one batch of CRV,
directed, or formal work at a time. The next claim is a large coverage space
and a large formal partition under heterogeneous parallelism---many simulator
slots, scarce formal licenses, and build/writer seats. CHIA's role is
admission and replanning: keep high-yield simulation busy while the formal seat
works the highest expected-closure batch, cancel stale jobs when debt closes,
and reuse cached results keyed by immutable job identity. A timeout or missing
monitor must remain an open obligation, not a silent demotion to a narrower
scope under the same result label.

\section{Evidence and scope discipline}

Bounded cover non-hits do not retire coverage. Exclusion requires a separately
reviewed proof of the exact obligation under matching scope and assumptions.
\texttt{BOOM\_FORMAL\_SCOPE=lsu} is an explicit lower-scope experiment; optional
timeout fallback requires \texttt{BOOM\_FORMAL\_ALLOW\_LSU\_FALLBACK=1} and
cannot replace a ChipTop result. Historic 51-point hit/closure figures from the
proxy meter are retained for inspectability and are not submission ROI.

\section{BOOM collateral and qualification status}

The package is pinned to BOOM revision
\texttt{58ef2720eae13be26b3008c02b5a74ce29c61c44}. It contains 67 assertion
families, 66 scenario families, and activation identities, with overlay
staging and an explicit last-mile obligation list. For the pinned
MediumBoomV3Config, Chisel elaboration and five assertion-enabled,
coverage-instrumented DUT workloads completed: two directed programs, one
\texttt{riscv-dv} program, and two scalar-FP torture programs. The merged full-design
signal-toggle rate is 73.59\%; a filtered 15-module monitor subset reaches
77.96\%. These are activity metrics, not semantic property coverage. There
are 582 hits among 1,192 named cover counters, including 49/67 activation
families and 24/66 scenario families. The FP BusyTable activation reaches
63/64 instances for each family; the preg-0 instance remains unhit and is not
retired.

Local DTLB cover reached all nine lowered cover goals and the late-refill
witness replayed. A one-expression ROB rollback mutant triggered its intended
checker on the same directed workload that passes on baseline. Both generated
integer and FP RenameBusyTable assertion cones pass Yosys 2-inductive proofs.
These are partial, targeted results, not whole-processor qualification.

\begin{itemize}
\raggedright
\item \textbf{Elaboration and simulation:} MediumBoomV3Config elaborated; five
  workloads passed. Merged signal-toggle coverage is 73.59\% full-design and
  77.96\% for the filtered 15-module monitor subset. These are activity
  measurements, not property coverage. The aggregate hits 582/1,192 named
  cover counters (49/67 activation families; 24/66 scenario families).
\item \textbf{DTLB witness:} all nine lowered cover goals were reached, and
  the late-refill witness replays. A depth-5 BMC counterexample to
  \texttt{no\_stale\_refill\_after\_fence} also replays and needs owner review.
\item \textbf{ROB mutation:} one internal rollback mutant triggers the
  unchanged \texttt{rob.rollback\_no\_commit} checker; the same workload passes
  on baseline.
\item \textbf{Formal limits:} a depth-2 BoomFrontend assertion BMC and depth-6
  BusyTable BMCs pass. Both generated BusyTable assertion cones also pass
  2-inductive proofs, but no whole-core or DTLB induction proof has run.
  Per-instance checker qualification remains incomplete; the DTLB stale-refill
  counterexample needs BOOM-owner protocol review.
\item \textbf{Campaign limits:} the paired B1/B2/B3 compute-hour campaign on
  $N$ workers has not run.
\end{itemize}

\section{Reproduction}

The exact targeted Python regression command is in the repository's
\texttt{QUALIFICATION.md}; the Yosys BusyTable proof commands and logs are in
\texttt{formal\_busytable/README.md}. The test suite exercises package and
dispatch plumbing, not the N-worker campaign. The full adaptive-versus-static
protocol, including cold/warm scenarios and mechanism ablations, is specified
in the companion whitepaper in the artifact. BOOM qualification logs, hashes,
replay traces, and reproduction instructions are in
\texttt{artifacts\_mediumboom/kapil\_qualification/}.

\section{Conclusion}

We restate the hackathon contribution as (1)~a fair three-arm
coverage-per-compute-hour design in which CHIA must beat a fixed hybrid that
already has formal, and (2)~a path from single-seat serial placement to
parallel multi-resource CHIA on large cover and formal work. The BOOM monitor
package is submitted as reviewable next-phase collateral under an explicit
qualification boundary.

\section*{AI assistance}
Drafting and copy-editing used coding assistants. The authors are responsible
for the claims and evidence boundaries in this paper.

\end{document}
